\documentclass[11pt]{article}

\usepackage{amsmath,amsthm,amssymb}
\usepackage[margin=1.2in]{geometry}
\usepackage{hyperref}

\newtheorem{theorem}{Theorem}
\newtheorem{proposition}[theorem]{Proposition}
\newtheorem{corollary}[theorem]{Corollary}
\newtheorem{lemma}[theorem]{Lemma}
\newtheorem{conjecture}[theorem]{Conjecture}
\theoremstyle{definition}
\newtheorem{definition}[theorem]{Definition}
\theoremstyle{remark}
\newtheorem{remark}[theorem]{Remark}

\title{The Total Rank Formula and Image Basis Conjecture\\for the Staircase Product}
\author{Clio}
\date{April 16, 2026}

\begin{document}
\maketitle

\begin{abstract}
We reduce the $H$-invariant theorem
$\operatorname{rank}(\Pi|_{V_\lambda}) = \dim(V_\lambda^H)$
to a single statement: the \emph{image basis conjecture},
which asserts that right multiplication by the staircase product
$\Pi \in \mathbb{Q}[S_n]$ sends the descent-paired basis elements
to linearly independent vectors.
Combined with the eigenvalue positivity theorem (proved via
Kazhdan--Lusztig positivity), this yields the full theorem.
The image basis conjecture is verified computationally for $n \leq 7$.
\end{abstract}

\section{Setup and statement}

Let $\Pi = \prod_{k=1}^{n-1} \prod_{j=k}^{1} (1 + s_j) \in \mathbb{Q}[S_n]$
be the staircase product, where $s_j = (j, j\!+\!1)$ is the simple transposition.
Let $H = \langle s_1, s_3, s_5, \ldots \rangle \cong (\mathbb{Z}/2)^{\lfloor n/2 \rfloor}$.

\begin{definition}
A permutation $w \in S_n$ is \emph{descent-paired} if $w(2i-1) > w(2i)$
for all $1 \leq i \leq \lfloor n/2 \rfloor$.
Let $D \subset S_n$ denote the set of descent-paired permutations.
\end{definition}

\begin{proposition}\label{prop:Dprops}
$|D| = n!/2^{\lfloor n/2 \rfloor}$, and $D$ is a complete set of
representatives for the right cosets of~$H$ in~$S_n$.
\end{proposition}

\begin{conjecture}[Image Basis Conjecture]\label{conj:image}
The vectors $\{e_w \cdot \Pi : w \in D\}$ are linearly independent
in~$\mathbb{Q}[S_n]$.
\end{conjecture}

\begin{theorem}[Total Rank Formula]\label{thm:totalrank}
If Conjecture~\ref{conj:image} holds for~$n$, then
\[
  \operatorname{rank}(\Pi \text{ on } \mathbb{Q}[S_n]) = \frac{n!}{2^{\lfloor n/2 \rfloor}}.
\]
\end{theorem}

\begin{theorem}[$H$-Invariant Theorem]\label{thm:Hinvariant}
Assuming Conjecture~\ref{conj:image} for~$n$: for every partition
$\lambda \vdash n$,
\[
  \operatorname{rank}(\Pi|_{V_\lambda}) = \dim(V_\lambda^H).
\]
\end{theorem}

\section{Proofs}

\subsection{Positive semidefiniteness of the convolution form}

Write $\Pi = \sum_{g \in S_n} f_g \, g$ where $f_g > 0$ for all~$g$
(since every permutation appears with positive coefficient in the
staircase product, being the evaluation of a Schubert polynomial at~$1$).

\begin{lemma}\label{lem:psd}
The bilinear form $B(x, y) = \langle x, R_\Pi \, y \rangle$
on $\mathbb{Q}[S_n]$ is positive semidefinite, where $R_\Pi$
denotes right multiplication by~$\Pi$ and $\langle \cdot, \cdot \rangle$
is the standard inner product on the group algebra
(with $\{e_w\}$ orthonormal).
\end{lemma}

\begin{proof}
The right multiplication operator $R_\Pi$ on the regular representation
decomposes as $\bigoplus_\lambda \mathrm{Id}_{V_\lambda} \otimes \rho_{\lambda^*}(\Pi)$,
where $\lambda^*$ is the conjugate partition.
By the eigenvalue positivity theorem~\cite{Clio-EP},
all nonzero eigenvalues of $\rho_\mu(\Pi)$ are positive for every~$\mu$.
Hence $R_\Pi$ is positive semidefinite.
\end{proof}

\begin{lemma}\label{lem:symmetric}
$f_g = f_{g^{-1}}$ for all $g \in S_n$.
Hence the matrix $M_R[v,w] = f_{w^{-1}v}$ is symmetric.
\end{lemma}

\begin{proof}
The staircase reduced word $\mathbf{i} = (i_1, \ldots, i_m)$ has the
property that its reversal $(i_m, \ldots, i_1)$ is also a reduced word
for~$w_0$.
A subset $J \subseteq \{1, \ldots, m\}$ gives a reduced expression
for~$w$ in the forward reading if and only if $\{m+1-j : j \in J\}$
gives a reduced expression for~$w^{-1}$ in the reverse reading.
Since $\Pi$ counts subwords of the staircase word,
$f_g = f_{g^{-1}}$.
(Alternatively, this is verified computationally for all $n \leq 7$.)
\end{proof}

\subsection{Reduction to the image basis conjecture}

\begin{proof}[Proof of Theorem~\ref{thm:totalrank}]
By Lemma~\ref{lem:psd}, $R_\Pi$ is positive semidefinite.
The submatrix $M_R[D,D]$ with entries $f_{w^{-1}v}$ for $v, w \in D$
is a principal submatrix of a PSD matrix, hence PSD.

Conjecture~\ref{conj:image} states that the columns of the
right multiplication matrix $R_\Pi$ indexed by~$D$ are linearly independent.
This holds if and only if $M_R[D,D]$ is positive definite
(since for a PSD form, the Gram matrix of a subset of vectors is PD
iff the vectors are linearly independent).

If $M_R[D,D]$ is PD, then $\operatorname{rank}(R_\Pi) \geq |D| = n!/2^{\lfloor n/2 \rfloor}$.

For the upper bound: $\Pi \cdot s_1 = \Pi$ (the rightmost factor of
each block is $(1+s_1)$, which absorbs $s_1$).
Hence $R_\Pi(e_w - e_{s_1 w}) = 0$ for all~$w$,
giving $\operatorname{rank}(R_\Pi) \leq n!/2$.
More generally, the cascade analysis~\cite{Clio-cascade} yields
the per-irrep upper bound
$\operatorname{rank}(\Pi|_{V_\lambda}) \leq \dim(V_\lambda^H)$
via the eigenvalue positivity theorem.
Summing over irreducibles using the regular representation
decomposition:
\[
  \operatorname{rank}(R_\Pi)
  = \sum_\lambda \dim(V_\lambda) \cdot \operatorname{rank}(\Pi|_{V_\lambda})
  \leq \sum_\lambda \dim(V_\lambda) \cdot \dim(V_\lambda^H)
  = \frac{n!}{|H|}
  = \frac{n!}{2^{\lfloor n/2 \rfloor}}.
\]
The last equality is Frobenius reciprocity applied to
$\operatorname{Ind}_H^{S_n}(\mathbf{1})$.
\end{proof}

\begin{proof}[Proof of Theorem~\ref{thm:Hinvariant}]
By Theorem~\ref{thm:totalrank}:
\[
  \sum_\lambda \dim(V_\lambda) \cdot \operatorname{rank}(\Pi|_{V_\lambda})
  = \frac{n!}{2^{\lfloor n/2 \rfloor}}
  = \sum_\lambda \dim(V_\lambda) \cdot \dim(V_\lambda^H).
\]
Since $\operatorname{rank}(\Pi|_{V_\lambda}) \leq \dim(V_\lambda^H)$
for every~$\lambda$ (eigenvalue positivity + upper bound comparison~\cite{Clio-EP}),
equality of the weighted sums forces
$\operatorname{rank}(\Pi|_{V_\lambda}) = \dim(V_\lambda^H)$
for every~$\lambda$.
\end{proof}

\section{Computational verification}

The image basis conjecture is equivalent to positive definiteness of
$M_R[D,D]$.
We verify this by computing exact determinants:

\begin{center}
\begin{tabular}{cccccc}
\hline
$n$ & $|D|$ & $\det(M_R[D,D])$ & Factored & Min eigenvalue & PD? \\
\hline
3 & 3 & 4 & $2^2$ & 1.000 & Yes \\
4 & 6 & 576 & $2^6 \cdot 3^2$ & 0.658 & Yes \\
5 & 30 & $2^{70} \cdot 3^{10}$ & --- & 0.693 & Yes \\
6 & 90 & $2^{272} \cdot 3^{52} \cdot 5^8 \cdot 7^2 \cdot 11^8$ & --- & 0.551 & Yes \\
7 & 630 & $> 10^{300}$ & --- & 0.584 & Yes \\
\hline
\end{tabular}
\end{center}

\noindent
Key structural properties of $M_R[D,D]$:
\begin{itemize}
\item Symmetric ($f_g = f_{g^{-1}}$).
\item Constant diagonal: $M_R[w,w] = f_e$ for all $w \in D$.
\item Positive definite for all tested~$n$.
\item The minimum eigenvalue appears to remain bounded away from~$0$
(between $0.5$ and~$1.0$ for all $n \leq 7$).
\end{itemize}

\section{Discussion}

\paragraph{Why right multiplication.}
The left multiplication submatrix $M_L[D,D]$ (with entries $f_{vw^{-1}}$)
is singular for $n \geq 3$: distinct descent-paired permutations can
yield identical rows.
The right multiplication submatrix $M_R[D,D]$ (with entries $f_{w^{-1}v}$)
avoids this degeneracy and is positive definite.
The crucial asymmetry arises because $f_{vw^{-1}} \neq f_{w^{-1}v}$
in a non-abelian group.

\paragraph{Structure of the determinant.}
The determinants factor into powers of small primes, suggesting a
product formula.
The diagonal value $f_e$ (the coefficient of the identity in~$\Pi$,
counting subwords of the staircase giving~$e$) satisfies
$f_e = 2, 6, 34, 342, 6408$ for $n = 3, 4, 5, 6, 7$.

\paragraph{Proof status.}
The full $H$-invariant theorem (Theorem~\ref{thm:Hinvariant}) is
\textbf{proved} modulo Conjecture~\ref{conj:image}.
All other ingredients---eigenvalue positivity, the upper bound,
the Frobenius squeeze---are established unconditionally.
Conjecture~\ref{conj:image} is verified through $n = 7$
($|D| = 630$, exact integer determinant nonzero).

\paragraph{Possible proof strategies for Conjecture~\ref{conj:image}.}
\begin{enumerate}
\item \textbf{Cascade induction}: Track linear independence of
$\{e_w \cdot \Pi_{\leq k} : w \in D\}$ through the staircase factors,
showing independence is preserved at each step.
\item \textbf{Total positivity}: The factored structure of~$\Pi$
as a product of elementary ``row addition'' maps may yield total
nonnegativity of~$M_R[D,D]$, implying positive definiteness.
\item \textbf{Representation-theoretic}: Use the Kazhdan--Lusztig
cell structure of the Hecke algebra to control the rank of~$\Pi$
on parabolic quotients.
\end{enumerate}

\begin{thebibliography}{99}

\bibitem[Clio-EP]{Clio-EP}
Clio, \emph{Eigenvalue positivity for the Hecke staircase element
via Kazhdan--Lusztig positivity}, April 2026.

\bibitem[Clio-cascade]{Clio-cascade}
Clio, \emph{Cascade proof of the $H$-invariant upper bound},
April 2026.

\end{thebibliography}

\end{document}
